\section{Related Work and State-of-the-Art Synthesis}
\label{sec:related_work}

\subsection{Cognitive Authentication and Shoulder-Surfing Defenses}
Authentication schemes designed to resist visual eavesdropping and screen recording traditionally rely on graphical passcodes, randomized grid matrices, or cognitive challenge-response transformations~\cite{abdelkarim2024online, binbeshr2025challenge, duzgun2022shoulder}. While early graphical protocols reduce exposure to physical observers, their deployment in production web environments has encountered persistent barriers. Entry completion times in multi-round graphical schemes frequently exceed fifteen seconds, memorability degrades over inactive intervals, and repeated ceremonies remain vulnerable to multi-session recording intersection attacks when challenge entropy or symbol multiplicity is low~\cite{duzgun2022shoulder, binbeshr2025challenge}.

Subsequent challenge-response proposals introduced mental arithmetic and visual permutation matrices where the authentication server renders a randomized symbol grid and the user executes an internal cognitive mapping, such as modular offset addition or geometric path traversal, to derive one-time response characters~\cite{binbeshr2025challenge}. However, prior cognitive schemes remained isolated standalone prototypes. They neglected integration with enterprise web federation protocols, lacked cryptographic token binding to protect post-authentication sessions, and remained exposed to automated intersection convergence when anchor references were publicly visible on screen. Even when architectures attempt to establish isolated browser authentication channels for passwords and FIDO2~\cite{gautam2025passwords} or deploy hardware-anchored FIDO2 protocols with provable cryptographic bounds~\cite{bindel2023fido2}, empirical evaluations by Kuchhal et al. and Mahfouz et al. uncover persistent relying party implementation flaws and passkey interoperability barriers~\cite{kuchhal2023evaluating, mahfouz2025passkeys}. Furthermore, Mahdad et al. demonstrate that unprivileged malware on client devices can execute UI overlay deception attacks against authenticators with limited displays without requiring root privileges~\cite{mahdad2024breaching}. Related approaches in FIDO2 account recovery leverage electronic identity tokens (eIDs) to avoid insecure fallback channels, confirming that cryptographic primitives require careful user-facing mediation~\cite{schwarz2022feido}. Table~\ref{tab:related_cognitive} synthesizes how BC-MFA addresses gaps left by prior cognitive and browser-channel defenses.

\begin{table}[!t]
\caption{Synthesis of Prior Cognitive and Browser-Channel Defenses vs. BC-MFA}
\label{tab:related_cognitive}
\centering
\resizebox{\columnwidth}{!}{%
\begin{tabular}{@{}lcccc@{}}
\toprule
\textbf{Approach} & \textbf{Est. Time} & \textbf{Intersection Def.} & \textbf{OIDC / DPoP} & \textbf{Zero HW} \\
\midrule
Banking Graphical~\cite{abdelkarim2024online} & $>12\text{ s}$ & Low (Static grid) & No & Yes \\
AR Shoulder-Surf.~\cite{duzgun2022shoulder} & 8--15 s & Medium (AR optics) & No & No (AR) \\
Grid PIN Ch.-Resp.~\cite{binbeshr2025challenge} & 4--7 s & Medium (Public pos.) & No & Yes \\
Browser Auth Ch.~\cite{gautam2025passwords} & $<3\text{ s}$ & N/A (Browser mod) & Partial & Partial \\
\textbf{BC-MFA (Ours)} & \textbf{3.3--5.0 s*} & \textbf{High ($d{=}4, L{=}4, k_{\max}$)} & \textbf{Yes (RFC 9449)} & \textbf{Yes} \\
\bottomrule
\multicolumn{5}{@{}l}{\footnotesize *Analytical literature-derived latency envelope; empirical study in Section~\ref{sec:experimental_design_and_thesis}.}
\end{tabular}%
}
\end{table}

\subsection{Second-Factor Weaknesses on Compromised Hosts}
Multi-Factor Authentication deployed on commodity endpoints faces severe threats from resident malware and adversary-in-the-middle toolkits. Large-scale empirical measurements of SMS phishing infrastructure by Nahapetyan et al. and supply-chain analyses of commodity phishing kits by Tejaswi et al., Lee et al., and Lim et al. demonstrate that telephonic and one-time factors collapse systematically under automated reverse-proxy campaigns and carrier-level SIM swapping~\cite{nahapetyan2024sms, tejaswi2022leaky, lee2024beneath, lim2025phishers}. Push-based mobile authenticators suffer from notification fatigue and concurrent prompt abuse, where attackers trigger parallel verification requests that users approve inadvertently~\cite{jubur2021bypassing, mahdad2023breaking}. Even enhanced variations incorporating compare-and-confirm number matching introduce cognitive friction that adversaries exploit through delayed relay prompts~\cite{jubur2025usability}. Additionally, Amft et al. demonstrate that enterprise account recovery workflows routinely reopen insecure fallback channels, such as SMS OTP or security questions, undoing the protection established by primary factors~\cite{amft2023weve}. Comprehensive evaluations by Jubur et al. confirm that legacy two-factor authenticators and password managers lack host-level isolation when the local operating environment is hostile~\cite{jubur2025indepth}.

\subsection{Behavioral Biometrics and Continuous Verification}
Behavioral biometrics capture neuro-muscular interaction patterns during user engagement, providing passive identity and liveness signals~\cite{syed2022zero}. Within web and mobile environments, dominant modalities comprise keystroke kinetics, mouse dynamics, and touchscreen gestures~\cite{incel2021dakota, khan2024mouse}. Keystroke analysis examines key-down duration (hold time) and inter-key latency (flight time), whereas pointer and touch tracking measure trajectory curvature, velocity inflection points, and capacitive contact area~\cite{incel2021dakota, khan2024mouse}.

Cognitive authentication ceremonies introduce distinct neuro-cognitive liveness markers. The human mental computation phase produces observable hesitation intervals and cursor deceleration immediately preceding input execution, followed by short keystroke bursts~\cite{khan2024mouse}. Automated infostealers, headless browser scripts, and replay bots struggle to reproduce this coupled visual-inspection and motor cadence simultaneously: scripted replays exhibit either near-zero hesitation or synthetic uniform delays detectable by anomaly classifiers~\cite{incel2021dakota, khan2024mouse}. However, because prior surveys evaluate continuous multi-minute sessions rather than short 4-character inputs on unfamiliar hardware, BC-MFA deploys behavioral telemetry primarily for human-vs-automation and relay anomaly screening rather than standalone cross-device biometric identification.

\subsection{Web Identity Federation and Protocol Standards}
Enterprise web and mobile environments depend on OpenID Connect Core 1.0 and OAuth 2.0 to govern delegated authentication between Identity Providers and Relying Parties~\cite{jannett2024sok, oidc_core1, rfc6749oauth}. OpenID Connect standardizes ID Token claims, IETF RFC 7636 mandates Proof Key for Code Exchange (PKCE) for public clients~\cite{rfc7636pkce}, and IETF RFC 8176 establishes the IANA Authentication Method Reference (\texttt{amr}) registry~\cite{rfc8176amr}. Automated compliance evaluations by Philippaerts et al. and formal channel-binding analyses by Golaszewski et al. indicate widespread token handling oversights in commercial OAuth and FIDO deployments, reinforcing the necessity of strict cryptographic binding~\cite{philippaerts2022oauch, golaszewski2026binding}. Furthermore, Hilbig et al. demonstrate that without Continuous Access Evaluation or proof-of-possession token binding, stolen bearer tokens grant persistent unauthorized access after the initial login completes~\cite{hilbig2023protect}.

When users connect from unmanaged endpoints, OAuth 2.0 Security Best Current Practice (RFC 9700) recommends step-up authorization, mandatory PKCE, and sender-constrained tokens~\cite{rfc9700bcp}. BC-MFA defines an OIDC public-client integration profile combining registered RFC 8176 \texttt{amr} claims with RFC 7636 PKCE and RFC 9449 Demonstrating Proof-of-Possession (DPoP) access token binding~\cite{rfc7636pkce, rfc8176amr, rfc9449dpop}, allowing relying parties to verify step-up assurance without requiring client-side cryptographic hardware.
